% Generated draft: requires mathematical and source review.
\documentclass[11pt]{article}
\usepackage[margin=25mm]{geometry}
\usepackage{fontspec}
\setmainfont{Latin Modern Roman}
\usepackage{amsmath,amssymb,mathtools,unicode-math}
\setmathfont{Latin Modern Math}
\usepackage{enumitem,needspace,xurl,hyperref}
\setlength{\parindent}{0pt}
\setlength{\parskip}{5pt}
\setlength{\emergencystretch}{3em}
\newcommand{\N}{\mathbb{N}}
\newcommand{\Z}{\mathbb{Z}}
\newcommand{\Q}{\mathbb{Q}}
\newcommand{\R}{\mathbb{R}}
\newcommand{\C}{\mathbb{C}}
\newcommand{\F}{\mathbb{F}}
\newcommand{\A}{\mathbb{A}}
\newcommand{\PP}{\mathbb{P}}
\newcommand{\E}{\mathbb{E}}
\newcommand{\eps}{\varepsilon}
\newcommand{\dd}{\,\mathrm{d}}
\begin{document}
\section*{UnsolvedMath ID: 30001410}
\textbf{Source problem code:} catalogue entry
\subsection{Definitions and mathematical statement}
Let \(C\subset\R^n\) be a convex body, understood in the usual sense of a compact convex set with nonempty interior, and put \(\Omega=\operatorname{int}C\). The supplied record does not specify the range of dimensions \(n\) or impose boundary regularity. For distinct \(x,y\in\Omega\), let \(a,b\in\partial C\) be the endpoints of the chord through them, ordered \(a,x,y,b\). With the conventional normalization, the Hilbert distance is
\[
 d_C(x,y)=\frac12\log\frac{\lVert y-a\rVert\,\lVert x-b\rVert}{\lVert x-a\rVert\,\lVert y-b\rVert}.
\]
For \(x\in\Omega\) and \(v\ne0\), let \(t_+,t_->0\) satisfy \(x+t_+v,x-t_-v\in\partial C\). The corresponding Hilbert norm is
\[
 F_C(x,v)=\frac12\left(t_+^{-1}+t_-^{-1}\right),\qquad F_C(x,0)=0.
\]
Under the usual induced one-dimensional Holmes--Thompson convention, the length of a piecewise \(C^1\) curve \(\gamma:[\alpha,\beta]\to\Omega\) is
\[
 \operatorname{vol}_1(\gamma)=\int_\alpha^\beta F_C\bigl(\gamma(t),\dot\gamma(t)\bigr)\,\dd t.
\]
The source does not define the carrier or admissibility conditions of its Crofton measure. A conventional affine-hyperplane interpretation can be made explicit as follows. Let \(\mathcal H(\Omega)\) be the space of unoriented affine hyperplanes meeting \(\Omega\), equipped with its usual Borel structure. For a curve \(\gamma\), put
\[
 N_\gamma(H)=\#\{t\in(\alpha,\beta):\gamma(t)\in H\},
\]
counting parameter occurrences and assigning \(+\infty\) when the set is infinite. On this interpretation, a positive Crofton representation means a nonnegative Borel measure \(\mu\) such that
\[
 \operatorname{vol}_1(\gamma)=\int_{\mathcal H(\Omega)}N_\gamma(H)\,\dd\mu(H)
\]
for every piecewise \(C^1\) curve in \(\Omega\), with any fixed positive normalization constant absorbed into \(\mu\).

\textbf{Source assertion to prove or disprove.} In every intended dimension, the existence of a positive Crofton measure for Holmes--Thompson curve length in the Hilbert geometry of \(C\) implies that \(C\) is an ellipsoid. Here an ellipsoid means a set
\[
 C=\{x\in\R^n:(x-c)^{\mathsf T}A(x-c)\leq1\},
\]
where \(c\in\R^n\) and \(A\) is a symmetric positive-definite real matrix. The source describes this conclusion equivalently as the Hilbert geometry being hyperbolic.

\textbf{Unresolved formulation qualification.} The hyperplane carrier, curve class, intersection convention, and measure conditions above specify a conventional candidate formulation; they are not supplied definitions from the record. Their agreement with the intended conjecture requires checking the original report. The asserted direction is necessity only. The record supplies neither a precise meaning of ``hyperbolic'' nor dimension restrictions for that equivalence.
\subsection{Short English statement}
The proposed conjecture says that a positive Crofton formula for Holmes--Thompson curve length in a Hilbert geometry forces the defining convex body to be an ellipsoid. The supplied record leaves the precise Crofton framework and dimension range unspecified, so the explicit hyperplane formulation above remains provisional. This is a statement-only draft requiring human mathematical review.
\subsection{Sources}
\begin{itemize}
\item Local UnsolvedMath registry: supplied complete record for permanent ID \(30001410\), title \emph{Crofton Measures in Hilbert Geometry}, source problem code \texttt{catalogue entry}, in the Oberwolfach Reports Open Problems collection. The proposed year is 2010; no individual proposer is identified.
\item Supplied report citation: \emph{Mini-Workshop: Valuations and Integral Geometry} (2010), Oberwolfach Reports, pp.~141--178, DOI 10.4171/owr/2010/04. Original Oberwolfach problem report: \url{https://doi.org/10.4171/owr/2010/04}. No specific problem page or passage is supplied.
\item Additional supplied literature source, labeled ``Schneider polytopal Hilbert geometries'': \url{https://home.mathematik.uni-freiburg.de/rschnei/Polytopal.Hilbert.pdf}. The record supplies no detailed theorem or bibliographic citation for this link.
\item Literature assessment checked at \texttt{2026-08-21}: the supplied dated curation labels the problem open, says that the report presents the ellipsoid-only assertion as a conjecture, and reports later existence results in special settings without a proof or counterexample in the stated generality. This is supplied dated curation, not an independently verified literature assessment; the record explicitly requests independent verification of time-sensitive status and cited claims.
\item The registry labels the statement \texttt{corrected\_verified}, explaining only that TeX notation was enclosed in mathematical delimiters without changing mathematical content. This label does not establish mathematical verification of the conjecture or the provisional definitions above.
\end{itemize}
\end{document}
